\chapter{Stratégies fondées sur des données alternatives}\label{ch:s1:alternative-data-strategies}

Lorsque les images satellite des parkings des enseignes de distribution sont parvenues pour la première fois aux investisseurs, ceux qui les achetaient pouvaient prendre position avant les rapports trimestriels des enseignes,
surtout lorsque ceux-ci étaient mauvais~; Katona, Painter, Patatoukas et Zeng ont relevé davantage de ventes à découvert informées et une liquidité plus faible autour des rapports des enseignes couvertes,
et en ont conclu qu'un accès inégal à ces données accroissait l'asymétrie d'information sans améliorer immédiatement la formation des prix. Un jeu de données alternatif
vaut le plus pour ses premiers acheteurs, et moins pour chacun de ceux qui suivent. Sur le marché synthétique, un panel qui lit un dixième des surprises de résultats des entreprises dix jours
à l'avance obtient un ratio de Sharpe de 4.4 lorsque personne d'autre ne le possède, de 1.9 lorsque la moitié du marché en dispose, et il perd de l'argent lorsque neuf dixièmes en disposent~; une firme qui reçoit les mêmes données un
jour avant tout le monde conserve un ratio de 4.3 quelle que soit la diffusion. Ce chapitre transforme un jeu de données en position et mesure comment sa valeur s'érode à mesure qu'il se diffuse. La
construction est \texttt{firm.altstrat}.

\section{De la donnée à la position}

\begin{definition}[Nowcast, stratégie fondée sur des données alternatives]\label{def:s1:alternative-data-strategies:nowcast}
Un \emph{nowcast}\index{nowcast} est une estimation d'une grandeur qui sera publiée plus tard (ventes, résultats, activité économique), établie à partir de données observées
avant le rapport. Une \emph{stratégie fondée sur des données alternatives}\index{stratégie fondée sur des données alternatives} prend position sur un titre à partir de nowcasts construits à partir de données extérieures aux
sources habituelles des marchés et de la comptabilité (transactions, images, fréquentation de sites web, textes), avant le rapport qui révélera la grandeur.
\end{definition}

Le chemin qui mène d'un jeu de données à une position comporte quatre étapes. Associer les données aux entreprises qu'elles mesurent (le livre~7, chapitre~12, présente la couverture et le
piège du point dans le temps que constituent les historiques reconstitués rétroactivement). Transformer les lectures en \omterm{def:s1:alternative-data-strategies:nowcast}{nowcast} de la grandeur publiée~: ici, la surprise de résultats. Transformer le \omterm{def:s1:alternative-data-strategies:nowcast}{nowcast} en
rendement espéré~: la réaction du cours à l'annonce par unité de surprise. Et tenir la position pendant la fenêtre au cours de laquelle le marché apprend ce que les données
disent déjà. Le panel du chapitre couvre un dixième des entreprises synthétiques (environ 420 annonces par an), lit chaque surprise avec un bruit égal au double de son écart-type (une
corrélation de 0.46), et est livré dix jours de négociation avant l'annonce. Le \omterm{def:s1:alternative-data-strategies:nowcast}{nowcast} est la lecture multipliée par une pente de 0.21 estimée sur les deux premières
années~; le portefeuille détient le \omterm{def:s1:alternative-data-strategies:nowcast}{nowcast} de la clôture de la livraison à la clôture de l'annonce, mis à l'échelle pour une exposition brute égale à un.

\section{Ce qui a été documenté comme efficace}

Les publications forment un ensemble de jeux de données, chacun lié à une grandeur que les cours n'intègrent que lentement. La couverture satellite des enseignes de distribution a permis à des investisseurs sophistiqués
de cibler les mauvais rapports (Katona et coauteurs). Les avis des salariés publiés sur des plateformes participatives permettent de prévoir des surprises de résultats~: Green, Huang, Wen et Zhou ont établi que les firmes dont
les notes se sont nettement améliorées surperformaient celles dont les notes avaient baissé, l'effet portant sur les avis des salariés actuels concernant les perspectives de carrière et la
direction générale. Les liens économiques en fournissent un exemple plus ancien~: Cohen et Frazzini ont constaté que les cours n'intègrent pas rapidement les informations concernant les principaux clients d'une firme,
et qu'une stratégie prenant position sur les fournisseurs en fonction des rendements de leurs clients a rapporté des alphas mensuels supérieurs à 150 points de base~; les liens entre fournisseurs du marché synthétique
(le livre~7, chapitre~10) reproduisent cet effet. Dans chaque cas, l'avantage tient à l'écart entre le moment où les données sont connues et celui où le cours les intègre.

\section{L'érosion comme diffusion des données}

\begin{definition}[Érosion des données]\label{def:s1:alternative-data-strategies:decay}
L'\emph{érosion des données}\index{erosion@érosion des données} est la perte de valeur d'un jeu de données pour la négociation, à mesure que davantage d'investisseurs l'acquièrent et prennent position plus tôt, faisant bouger les cours avant le
rapport et ne laissant qu'une part moindre de la réaction à chaque utilisateur~; à la différence de la décroissance après publication d'une anomalie publiée, elle est pilotée par le nombre d'acheteurs des
données.
\end{definition}

Zhu a mis en évidence l'autre versant du phénomène~: l'introduction des données de transactions de consommateurs et des données satellite a accru l'informativité des cours en abaissant le coût
d'acquisition de l'information. Des cours plus informatifs sont bénéfiques pour les marchés et néfastes pour la valeur du jeu de données. Le chapitre introduit directement le mécanisme
(\cref{lst:s1:alternative-data-strategies:diffuse})~: une part $\phi$ du marché achète le même panel et déplace chaque cours le jour de la livraison de $\phi$ de la
réaction espérée de l'annonce, compte tenu de la lecture~; cette part est alors absente de la réaction à l'annonce.

\begin{center}
\small
\begin{tabular}{@{}lccccc@{}}
\toprule
part du marché disposant des données & 0 & 0.25 & 0.5 & 0.75 & 0.9\\
\midrule
ratio de Sharpe après coûts & 4.41 & 3.20 & 1.86 & 0.45 & $-0.43$\\
rendement annuel après coûts (brut 1) & 61.1~\% & 42.6~\% & 24.2~\% & 5.7~\% & $-5.4~\%$\\
\midrule
avec les données un jour avant les autres~: ratio de Sharpe & 4.28 & 4.40 & 4.43 & 4.37 & 4.29\\
\bottomrule
\end{tabular}
\end{center}

La valeur diminue presque linéairement avec la part du marché qui détient les données, et devient négative avant que tout le monde ne les possède, parce que le portefeuille paie
dix points de base par unité échangée, pour un volume à sens unique de 36 fois son exposition brute par an. Le calendrier constitue l'autre moitié du problème. Un même jeu de données livré à la firme un jour avant les autres vaut
un ratio de Sharpe de 4.3 à 4.4 à tous les niveaux de diffusion (\cref{fig:s1:alternative-data-strategies:decay})~: la firme prend position avant le mouvement des autres
et le conserve jusqu'à son terme. Sur le marché des jeux de données, c'est pourquoi l'exclusivité et la rapidité de livraison ont un prix, et pourquoi les premiers clients des fournisseurs gagnent ce que leurs
clients ultérieurs ne peuvent pas obtenir.

\begin{omfigure}
\begin{tikzpicture}
\begin{axis}[width=10cm,height=5.4cm,xlabel={\small part du marché disposant des mêmes données},ylabel={\small ratio de Sharpe après coûts},
  tick label style={font=\footnotesize},xmin=0,xmax=0.95,ymin=-1,ymax=5,grid=major,grid style={gray!25},legend pos=south west,
  legend style={font=\scriptsize},table/col sep=comma]
\addplot[omThm,very thick,mark=*] table[x=phi,y=sr_net]{figdata/strategies-1/16-alternative-data-strategies/decay.csv};
\addplot[omDef,very thick,mark=square*] table[x=phi,y=sr_head]{figdata/strategies-1/16-alternative-data-strategies/decay.csv};
\addplot[black!40,dashed] coordinates {(0,0) (0.95,0)};
\legend{données le même jour que les autres,données un jour plus tôt}
\end{axis}
\end{tikzpicture}

\omcaption{Le portefeuille de \omterm{def:s1:alternative-data-strategies:nowcast}{nowcasts} sur le marché synthétique à mesure que le panel alternatif se diffuse~: ratio de Sharpe après coûts en fonction de la part du marché qui achète le
même panel, pour une firme qui le reçoit le même jour que les autres et pour une firme qui le reçoit un jour plus tôt. Données~:
\texttt{s1\_altstrat.run}.}\label{fig:s1:alternative-data-strategies:decay}
\end{omfigure}

\section{Exploiter une stratégie fondée sur des données alternatives}

Une \omterm{def:s1:alternative-data-strategies:nowcast}{stratégie fondée sur des données alternatives} en production relève autant de l'exploitation que du modèle. Les données arrivent en retard, changent de format, perdent leur couverture et sont révisées~; chaque livraison
doit être contrôlée par rapport à son propre historique avant d'affecter une position (le flux de travail du livre~7, chapitre~29). L'historique du fournisseur doit être établi au point dans le temps, ou le
backtest opère sur des lectures reconstituées rétroactivement et ajustées au résultat (le test du livre~7, chapitre~12). Le \omterm{def:s1:alternative-data-strategies:nowcast}{nowcast} doit être réestimé à mesure que les données changent. Et
l'érosion propre de la stratégie doit être suivie~: le modèle du chapitre suggère de surveiller la part de la réaction espérée qui intervient avant le rapport (le
mouvement d'avant-rapport par unité de \omterm{def:s1:alternative-data-strategies:nowcast}{nowcast}), laquelle croît à mesure que les données se diffusent et peut être observée sans savoir qui d'autre achète les données.

\section{Fiches de stratégie}

\begin{strategyfile}{Transaction de nowcast sur les revenus}\label{strat:s1:alternative-data-strategies:nowcast}
\sfield{Qui rémunère la stratégie, et pourquoi} Les investisseurs qui ne disposent pas des données et n'apprennent les résultats du trimestre qu'au moment du rapport.
\sfield{Instruments et plateformes} Actions des entreprises mesurées par le panel (enseignes de distribution, marques de grande consommation).
\sfield{Signal} Un \omterm{def:s1:alternative-data-strategies:nowcast}{nowcast} du chiffre d'affaires ou des résultats par rapport au consensus, établi à partir des transactions ou des images.
\sfield{Dimensionnement et exécution} Position de la livraison au rapport, proportionnelle au \omterm{def:s1:alternative-data-strategies:nowcast}{nowcast}~; couvrir le marché et le secteur.
\sfield{Coûts} Licences de données, souvent le poste le plus lourd~; négociation autour des rapports.
\sfield{Comment elle s'éteint} Diffusion des mêmes données à davantage d'acheteurs~; évolution de la couverture du panel.
\sfield{Horizon, capacité, infrastructure} Quelques semaines~; capacité limitée par les titres couverts~; une équipe d'ingénierie des données.
\sfield{Backtest honnête} Livraisons établies au point dans le temps, et non historique reconstitué~; dates des rapports telles qu'elles étaient programmées.
\sfield{Sources} Katona, Painter, Patatoukas et Zeng (2024)~; Zhu (2019)~; la simulation de ce chapitre.
\end{strategyfile}

\begin{strategyfile}{Signal de fréquentation web}\label{strat:s1:alternative-data-strategies:web}
\sfield{Qui rémunère la stratégie, et pourquoi} Comme pour les \omterm{def:s1:alternative-data-strategies:nowcast}{nowcasts}~: l'activité en ligne précède les ventes publiées.
\sfield{Instruments et plateformes} Entreprises dont les ventes se réalisent en ligne.
\sfield{Signal} Visites, usage des applications ou recherches, rapportés à leur profil saisonnier.
\sfield{Dimensionnement et exécution} Comme pour les \omterm{def:s1:alternative-data-strategies:nowcast}{nowcasts}.
\sfield{Coûts} Licences de données~; correspondance bruitée entre activité et chiffre d'affaires.
\sfield{Comment elle s'éteint} Changements dans la mesure de la fréquentation~; diffusion.
\sfield{Horizon, capacité, infrastructure} Quelques semaines~; mise en correspondance des domaines et des applications avec les entreprises.
\sfield{Backtest honnête} Fréquentation telle qu'elle était mesurée à l'époque, les changements de méthodologie du fournisseur étant signalés.
\sfield{Sources} Aucun chiffre de performance vérifié~; Green et al.\ (2019) pour le cas voisin des avis de salariés.
\end{strategyfile}

\begin{strategyfile}{Signal de sentiment}\label{strat:s1:alternative-data-strategies:sentiment}
\sfield{Qui vous paie, et pourquoi} Des investisseurs lents à agréger des opinions dispersées sur une entreprise.
\sfield{Instruments et plateformes} Des actions suffisamment commentées dans les avis, les publications ou les articles.
\sfield{Signal} Variations des notations participatives ou du ton des textes (livre~7, chapitre~12).
\sfield{Dimensionnement et exécution} Tris mensuels, ou inclinaisons dans un portefeuille plus large.
\sfield{Coûts} Faible rotation du portefeuille.
\sfield{Comment la stratégie meurt} Avis manipulés~; diffusion.
\sfield{Horizon, capacité, infrastructure} Des mois~; traitement de textes.
\sfield{Backtest honnête} Avis horodatés à leur publication~; avis supprimés et modifiés pris en compte.
\sfield{Sources} Green, Huang, Wen et Zhou (2019)~: les employeurs dont la note s'améliore surperforment nettement ceux dont la note se dégrade.
\end{strategyfile}

\begin{strategyfile}{Signal de liens de la chaîne d'approvisionnement}\label{strat:s1:alternative-data-strategies:supply}
\sfield{Qui vous paie, et pourquoi} Des investisseurs contraints par leur attention, qui ne suivent pas les nouvelles à travers les liens économiques.
\sfield{Instruments et plateformes} Des fournisseurs et leurs principaux clients.
\sfield{Signal} Les rendements ou les nouvelles récents des clients, rapportés à leurs fournisseurs.
\sfield{Dimensionnement et exécution} Position longue sur les fournisseurs des clients porteurs de bonnes nouvelles, position courte sur les autres~; mensuel.
\sfield{Coûts} Rotation modérée.
\sfield{Comment la stratégie meurt} Liens mieux connus~; données sur les liens désormais largement disponibles.
\sfield{Horizon, capacité, infrastructure} Un mois~; une base de données de la chaîne d'approvisionnement, reconstituée telle qu'elle était connue à chaque date.
\sfield{Backtest honnête} Liens tels qu'ils étaient publiés à chaque date~; moment de publication des informations sur les clients.
\sfield{Sources} Cohen et Frazzini (2008)~: des alphas mensuels supérieurs à 150 points de base~; livre~7, chapitre~10.
\end{strategyfile}

\section{Tutoriel~: au moment où l'on vend}\label{tut:s1:alternative-data-strategies:tut}

\begin{tutorial}
\textbf{Objectif.} Transformer un panel synthétique en positions de \omterm{def:s1:alternative-data-strategies:nowcast}{nowcast} et mesurer comment leur valeur diminue à mesure que les données se diffusent, ainsi que ce que vaut une longueur d'avance.
\textbf{État final~:} le tableau et la \cref{fig:s1:alternative-data-strategies:decay}.
\begin{tutsteps}
\item \textbf{Diffusion et carnet}~: les autres intervenants négocient les données à leur livraison~; le carnet est conservé de la livraison jusqu'au rapport.
\omcode{code/firm/altstrat/firm_altstrat.py}{33}{54}{Diffusion des données et carnet précédant l'événement.}\label{lst:s1:alternative-data-strategies:diffuse}
\item \textbf{Exécuter} \texttt{nowcast\_quality()}, \texttt{run(phi)} pour cinq actions et \texttt{run(phi, 1)} pour la longueur d'avance, puis \texttt{fig\_altstrat.py}.
\end{tutsteps}
\textbf{Pour aller plus loin.} Faire croître au fil des années la part des acteurs disposant des données et mesurer le ratio de Sharpe glissant du carnet~; estimer la diffusion à partir du
mouvement précédant le rapport par unité de \omterm{def:s1:alternative-data-strategies:nowcast}{nowcast}~; combiner deux panels aux délais de livraison différents.
\end{tutorial}

\section{Construction~: stratégies fondées sur des données alternatives}\label{bld:s1:alternative-data-strategies:altstrat}

\begin{build}
\bfield{Rôle} Des \omterm{def:s1:alternative-data-strategies:nowcast}{nowcasts} à partir d'un panel bruité, la diffusion de ces mêmes données à travers le marché, et le carnet précédant le rapport.
\bfield{Interface} \texttt{measure(surprise, noise, rng)}, \texttt{nowcast\_slope(readings, surprises)},
\texttt{diffuse(R, events, jump, phi, lead, reading, shrink)}, \texttt{pre\_event\_book(events, nowcast, lead, T, N)}.
\bfield{Règles} Les pentes de \omterm{def:s1:alternative-data-strategies:nowcast}{nowcast} sont estimées sur les seuls événements antérieurs~; la diffusion conserve le mouvement total de chaque annonce~; les positions sont prises à la clôture de la livraison.
\bfield{Tests d'acceptation} \texttt{code/firm/altstrat/tests/}~: la pente d'une lecture bruitée~; la diffusion déplace la part prévue et conserve le total~; les pondérations
du carnet vérifiées à la main.
\bfield{Pour aller plus loin} Diffusion croissante au fil du temps~; fournisseurs dont les historiques sont révisés~; plusieurs panels aux délais différents.
\end{build}

\begin{omsources}
\item C.~Zhu, ``Big data as a governance mechanism'', \emph{Review of Financial Studies} 32(5), 2019.
\item Z.~Katona, M.~O.~Painter, P.~N.~Patatoukas et J.~Zeng, ``On the capital market consequences of big data: evidence from outer space'', \emph{Journal
of Financial and Quantitative Analysis} 60(2), 2025 (en ligne en 2024).
\item L.~Cohen et A.~Frazzini, ``Economic links and predictable returns'', \emph{Journal of Finance} 63(4), 2008.
\item T.~C.~Green, R.~Huang, Q.~Wen et D.~Zhou, ``Crowdsourced employer reviews and stock returns'', \emph{Journal of Financial Economics} 134(1), 2019.
\end{omsources}

\section{Exercices}

\begin{exercise}[$\star$]\label{exo:s1:alternative-data-strategies:1}
Un panel lit une surprise entachée d'un bruit égal à deux fois l'écart type de la surprise. Calculer la corrélation de la lecture avec la surprise, ainsi que la meilleure pente de \omterm{def:s1:alternative-data-strategies:nowcast}{nowcast}.
\end{exercise}

\begin{exercise}[$\star$]\label{exo:s1:alternative-data-strategies:2}
La volatilité spécifique quotidienne d'une action est de 1.6~\% et le saut lié à l'annonce vaut trois fois celle-ci par unité de surprise. Avec une pente de \omterm{def:s1:alternative-data-strategies:nowcast}{nowcast} de 0.2 et une lecture de 2,
calculer la réaction attendue par le \omterm{def:s1:alternative-data-strategies:nowcast}{nowcast}, puis la part de cette réaction qui se produit à la livraison si la moitié du marché dispose des données.
\end{exercise}

\begin{exercise}[$\star$]\label{exo:s1:alternative-data-strategies:3}
Pourquoi l'\omterm{def:s1:alternative-data-strategies:decay}{érosion des données} diffère-t-elle de la décroissance après publication d'une anomalie~?
\end{exercise}

\begin{exercise}[$\star\star$]\label{exo:s1:alternative-data-strategies:4}
Pourquoi le carnet perd-il de l'argent avant que l'ensemble du marché dispose des données~?
\end{exercise}

\begin{exercise}[$\star\star$]\label{exo:s1:alternative-data-strategies:5}
Pourquoi une longueur d'avance d'un jour vaut-elle presque toute la stratégie~?
\end{exercise}

\begin{exercise}[$\star\star$]\label{exo:s1:alternative-data-strategies:6}
Comment estimer, à partir des seuls cours, la proportion d'acteurs qui utilisent vos données~?
\end{exercise}

\begin{exercise}[$\star\star\star$]\label{exo:s1:alternative-data-strategies:7}
\emph{Programmation.} Exécuter \texttt{run(0.5, 0, 3.0)}, avec un bruit du panel égal à trois fois l'écart type de la surprise au lieu de deux. Que devient le ratio de Sharpe, et pourquoi~?
\end{exercise}

\begin{exercise}[$\star\star\star$]\label{exo:s1:alternative-data-strategies:8}
\emph{Trouver la faille.} «~La série historique de cinq ans du fournisseur montre que notre \omterm{def:s1:alternative-data-strategies:nowcast}{nowcast} aurait atteint un ratio de Sharpe de 3~; nous paierons la licence.~»
\end{exercise}

\section{Problème~: au moment où l'on vend}

\begin{problem}[{Problème du week-end --- la vie d'un jeu de données}]\label{pb:s1:alternative-data-strategies:1}
Le panel synthétique du chapitre et le registre public.

\medskip\noindent\textbf{Partie I --- Des données à la position.}
\begin{enumerate}
\item Définir un \omterm{def:s1:alternative-data-strategies:nowcast}{nowcast} et une \omterm{def:s1:alternative-data-strategies:nowcast}{stratégie fondée sur des données alternatives}.
\item Énumérer les quatre étapes allant d'un jeu de données à une position.
\item Décrire le panel du chapitre et son \omterm{def:s1:alternative-data-strategies:nowcast}{nowcast}.
\item Que détient le carnet, et à quel moment~?
\end{enumerate}

\medskip\noindent\textbf{Partie II --- Le registre.}
\begin{enumerate}[resume]
\item Que Katona et ses coauteurs ont-ils trouvé sur les données satellitaires~?
\item Que Green et ses coauteurs ont-ils trouvé sur les avis concernant les employeurs~?
\item Que Cohen et Frazzini ont-ils trouvé sur les liens économiques~?
\item Que Zhu a-t-il trouvé sur le contenu informatif des prix~?
\end{enumerate}

\medskip\noindent\textbf{Partie III --- Érosion.}
\begin{enumerate}[resume]
\item Définir l'\omterm{def:s1:alternative-data-strategies:decay}{érosion des données} et décrire la manière dont le chapitre l'introduit.
\item Donner les ratios de Sharpe en fonction de la part de capital qui utilise les données, à mesure que celle-ci croît.
\item Pourquoi la valeur décroît-elle presque linéairement~?
\item Quel effet produit une longueur d'avance~?
\end{enumerate}

\medskip\noindent\textbf{Partie IV --- Le verdict.}
\begin{enumerate}[resume]
\item Énoncer le \emph{résultat nommé}~: le ratio de Sharpe de la stratégie de \omterm{def:s1:alternative-data-strategies:nowcast}{nowcast} en fonction de la part de capital informé qui utilise les données à mesure qu'elle croît.
\item Comment suivre l'érosion d'un jeu de données~?
\item Quels risques opérationnels la stratégie comporte-t-elle~?
\item Comment évaluer l'historique d'un fournisseur~?
\item Quelle valeur une licence de données doit-elle avoir pour la société~?
\item Quelle fiche de stratégie s'érode le plus vite~?
\item Quel lien ces stratégies entretiennent-elles avec le chapitre~7~?
\item En une phrase~: que vend un jeu de données alternatives~?
\end{enumerate}
\end{problem}

\section{Questions d'entretien}

\begin{interviewq}[$\star$ \iqroles{chercheur}]\label{iq:s1:alternative-data-strategies:1}
Comment évalueriez-vous un nouveau jeu de données alternatives~?
\end{interviewq}

\begin{interviewq}[$\star\star$ \iqroles{chercheur}]\label{iq:s1:alternative-data-strategies:2}
Qu'est-ce que le biais de remplissage rétroactif (backfill bias) dans les données d'un fournisseur, et comment le détectez-vous~?
\end{interviewq}

\begin{interviewq}[$\star\star$ \iqroles{chercheur, trader}]\label{iq:s1:alternative-data-strategies:3}
Votre signal fondé sur des données alternatives s'est affaibli sur deux ans. Que peut-il se passer~?
\end{interviewq}

\begin{interviewq}[$\star\star$ \iqroles{développeur}]\label{iq:s1:alternative-data-strategies:4}
Concevriez-vous le pipeline qui transforme une livraison quotidienne du fournisseur en transaction le soir même~?
\end{interviewq}

\begin{interviewq}[$\star\star$ \iqroles{risque}]\label{iq:s1:alternative-data-strategies:5}
Quels sont les risques juridiques et de conformité liés aux données alternatives~?
\end{interviewq}

\begin{interviewq}[$\star\star\star$ \iqroles{chercheur}]\label{iq:s1:alternative-data-strategies:6}
Une part $\phi$ du marché négocie sur un signal à sa livraison et fait bouger le prix de $\phi$ de la réaction espérée. Calculez l'espérance du P\&L d'un trader qui
le reçoit à la livraison, et celle d'un trader qui le reçoit $k$ jours plus tôt et conserve sa position jusqu'à celle-ci.
\end{interviewq}
